%
%============= CHAPTER BREAK =================%
%
\chapter{Aelius Promotus (ca 138--193 CE)}%
\vspace*{3mm} %
\justifying %

\noindent{}According to the \textit{Oxford Reference} entry, %
Aelius Promotus was a “physician from Alexandria, probably belonged to the period between Hadrian and Pertinax (ad 138--93)”.\endnote{%
Available at \url{www.oxfordreference.com/display/10.1093/oi/authority.20110803095353509}.%
} %
This reference entry is outdated, though, %
in regard to its claim that his work \textit{On Potency} / Δυναμερόν is unedited, %
given that Crismani 2002 contains just such an edition, %
one now available to query in the TLG. %
The treatise \textit{On Poisonous Animals and Toxic Drugs} / %
\textit{De venenatis animalibus et venenis medicamentis} %
is a different matter. %

While the latter treatise has traditionally been attributed to Promotus, %
this attribution was dubiously based on two manuscripts that set the treatise between two different works of Promotus.\endnote{%
Ihm 1995; Salazar 1998.} %
Ihmʼs critical edition (1995) and related article (1996) refute Rohde (1873, 1901), %
who had claimed that \textit{De venenatis} was composed by the 1st--2nd century CE medical author Archigenes, %
when in fact Archigenes was only occasionally excerpted, %
along with several other medical authors who date no later than the 2nd century CE. %
Ihm does not conclusively disqualify Promotus as author of this treatise, %
but cautiously maintains that its author is unknown, %
that it was composed and occasionally edited between the 2nd and 6th century CE. %
Given the complex, evolving character of the work, we indicate the author of \textit{Poisonous Animals} as [Ps-]Promotus. %
\\[1mm] %

\noindent{}\textbf{ἄψινθος} / \textit{apsinthos} / “wormwood”|“absinthe”, %
as a feminine, second declension noun (LSJ ἄψινθος, ἡ), %
has its first possible occurrence in non-Christian texts in the TLG in one of three candidates: %
the medical authors [Ps-]Promotus and Aretaeus, and the grammarian Herodian. %
While Aretaeus and Herodian are likely more secure chronological origin points for the feminine second-declension noun, %
we are exploring this formal neologism candidate in the chapter on Promotus because the usage by [Ps-]Promotus %
bears the closest resemblance to the usage found in the NT. %
Indeed, in our opinion the Promotus-affiliated tradition partly inspired the editor(s) of Revelation at this point. %

Before exploring the earliest history of the feminine second-declension noun, %
we must take note that the neuter second declension ἀψίνθιον / \textit{apsinthion} %
is attested in Attic Greek works by Xenophon (\textit{Anab}. 1.5.1) %
the medical author Hippocrates (15 TLG hits), %
the botanist Theophrastus (7 TLG hits), and others. %
The neuter lemma remains consistent in the mid-1st century CE medical author Dioscorides Pedanius (65 TLG hits), %
the first Greek author in the TLG to attest the derivative lemma ἀψινθίτης / \textit{apsinthitēs} / %
“absinthine wine” and the Latin loanword \textit{apsinthium}. %
The consistent use of the neuter second-declension continues in several other medical authors roughly contemporaneous to NT texts, %
such as: %

\begin{spacing}{0.8}
\begin{small}
\begin{itemize}
\item Ps-Dioscorides (10 TLG hits, post 1st c. CE) %
\item Rufus (7 TLG hits, 1st--2nd c. CE) %
\item Soranus (7 TLG hits, 1st--2nd c. CE) %
\item Archigenes (8 TLG hits, 1st--2nd c. CE) %
\item Galen (120 TLG hits, mid-2nd c. CE) %
\item Ps-Galen (23 TLG hits, post 2nd c. CE) %
\end{itemize}
\end{small}
\end{spacing} \\[1.5mm] %

Among Greek papyri, only the neuter form appears (11 hits), %
and only two of these are prior to the 2nd century CE (p.herc. 182, p.herc. 1471). %
Among non-canonical (or ambiguously canonical) Christian writings, %
the \textit{Shepherd of Hermas} also only uses the neuter second-declension (3 hits, all in ch. 33).\clearpage %

The neuter lemma appears only once in the [Ps-]Promotus treatise \textit{On Poisonous Animals and Toxic Drugs}, %
specifically when the author/compiler quotes the Classical medic Epaenetus (\textit{venen}. 63). %
Elsewhere the unusual spelling ἀψίνθιν / \textit{apsinthin} appears (\textit{venen}. 53), %
a rare feminine third-declension form, or perhaps a misspelling. %

What may be one of the earliest uses of the feminine second-declension noun ἄψινθος / %
\textit{apsinthos} %
is found in the same work by [Ps-]Promotus, %
in the section “On Mosquitoes” / πρὸς κώνωπας. %
It comes as the last of eight (!) different methods prescribed %
for keeping mosquitoes away from oneʼs home and body. %
The fifth method sounds somewhat similar to the sizzle of todayʼs electric mosquito zappers; %
it entailed placing a fuming oil lamp inside a pot with many holes, designed to attract and kill the insects. %

\vspace*{3mm} 
 \begin{quote}ἄλλο· χύτραν πολύοπον πάνυ κατασκευάσας καὶ †ἐμβαλὼν† λύχνον λάβρως ἀνάπτων, %
ἐπιτίθει κατ’ αὐτοῦ τὴν χύτραν κατακέφαλα ἐναποκλείων τὸν λύχνον. %
καὶ πάντες ἐκεῖ διαδραμόντες φλέγονται τὰς πτέρυγας. ... %
ἄλλο· \textbf{ἀψίνθην} σὺν ἐλαίῳ ἀπόζεσον καὶ χρῖε καὶ οὐχ ἅψονταί σου, ἢ ἔλαιον ῥαφάνινον. (\textit{venen}. 45) \\[1mm] %
Another [method]: Thoroughly crafting a many-eyed clay pot, %
and putting inside a lamp, furiously kindling it, % 
put down over it the clay pot, head down, enclosing the lamp inside. %
And there all of them, coursing through it, will have their wings burned. ... %
 Another [method]: \textbf{absinthe} with oil, boil to a froth and anoint [yourself], %
 and they will not touch you, or radish oil [as a substitute].\end{quote} % 
 \vspace*{3mm}

The usage of absinthe to deal with mosquitoes was not original to [Ps-]Promotus. %
The 1st century CE Dioscorides Pedanius may be the first known author to prescribe “absinthe with oil” / %
ἀψίνθιον σὺν ἐλαίῳ / \textit{apsinthion sun elaiō} (\textit{Eupor}. 2.137.1) for this. %
In Philumenusʼ late 2nd, early 3rd century CE treatise by a similar name to that of [Ps-]Promotus, %
\textit{On Venemous Animals and on Their Remedies} / \textit{De venenatis animalibus eorumque remediis}, %
“absinthe with oil in a soaked poultice” / ἀψίνθιον ϲὺν ἐλαίῳ ῥαφανίνῳ [καὶ] ϲυγχρίου / %
\textit{apsinthion sun elaiō raphaninō [kai] sugchriou} is twice mentioned as a preventative remedy for mosquitoes (12.2--3). 

In a later section “On Pine-Thistle” / περὶ ἰξίας, %
[Ps-]Promotus uses the feminine first-declension nominative, “absinthe with wine” / ἀψίνθη μετ’ οἴνου / \textit{apsinthē met oinou} (\textit{venen}. 59), %
specifically as a remedy to alleviate the poisonous effects of those who ingested this toxic plant. %
And again, in the section “On Deadly Mushrooms”, %
[Ps-]Promotus also uses the first-declension feminine accusative ἀψίνθην / \textit{apsinthēn}, %
specifically as as a cure for mushroom poisoning (\textit{venen}. 67). %
The latter belongs to a list of remedies that Promotus attributes %
to Epaenetus, though in this case, given the broader diachronic patterns of declension and spelling, %
it is very likely the spelling was updated from neuter second-declension to feminine first-declension. %

In his two books \textit{On the Cure of Chronic Diseases} / \textit{De curatione diuturnorum morborum}, %
the mid-2nd century CE medical author Aretaeus twice uses the feminine second-declension of “absinthe”, %
both with the formula “juice of \textbf{absinthe}” / χυλοῦ|χυλῷ τῆϲ \textbf{ἀψίνθου} / \textit{chulou|chulō tēs aphinthou}, %
once as a remedy for stomach bile (1.5.6) %
and once as a remedy for stomach ulcers (1.13.4). %
Elsewhere, Aretaeus favors the traditional and far more common neuter second-declension (7 TLG hits), %
three times as a “powder of absinthe” / ἀψινθίου κόμην|κόμῃ / \textit{apsinthiou komēn}|\textit{komē} (2.3.13, 2.4.6, 2.6.4), %
twice as a standalone reference to a liquid ἀψίνθιον / \textit{apsinthion} / “absinthe” (2.6.2, 2.7.3), 
one as “absinthe of wine” / ἀψινθίου τοῦ οἴνου / \textit{apsinthiou tou oinou} (2.13.11), %
and once as “juice of absinthe” / ἀψινθίου χυλὸν / \textit{apsinthiou chulon} (1.1.25). %

The same second declension neuter lemma also appears in the authentic \textit{General Prosody} of Herodian, %
who notes the variant spelling:\endnote{
The city name spelled with the \textit{upsilon} is also noted %
by Herodian in his likely authentic work \textit{On Orthography} (\textit{orth}. \textit{GG} 3.2:480). %
Regarding the authenticity of \textit{Orthog}., see Dyck 1993:788--789 and Dickey 2014:337--338.} %

\vspace*{2mm} 
 \begin{quote}Ἄψυνθος πόλις Θρᾴκης. %
 ἔστι καὶ εἶδος φυτοῦ. %
 διχῶς δ’ ἡ γραφὴ καὶ διὰ τοῦ υ καὶ διὰ τοῦ ι. (\textit{Pros.cath.}, \textit{GG} 3.1:147) \\[1mm] %
  Apsunthos is a city in Thrace. %
  It is also a type of plant. %
  But it is written in two ways, with an \textit{upsilon} and with an \textit{iota}.\end{quote} % 
 \vspace*{2mm}

\noindent{}Herodianʼs writings only apparently contain these descriptions %
of the neuter %
and feminine second-declension spellings for the lemmas associated with the plant, %
not the actual words themselves. %

Lemmas associated with absinthe occur in only one place in the NT, %
specifically in canonical Revelation, also known as the Apocalypse of John. %
As the LSJ entry for ἀψίνθιον indicates, %
Revelation is unique for pairing a feminine second-declension noun with a clearly masculine definite article. %
At least according to the NA\textsuperscript{28} edition of the Greek NT,\endnote{Aland et al 2012:753.} %
no notable variant spelling appears here. %
This makes it stand out in the sweep of linguistic history, %
but also likely reflects the unrefined Greek style of its editor(s) (Rev 8:10--11): %

\vspace*{2mm} 
 \begin{quote}
Καὶ ὁ τρίτος ἄγγελος ἐσάλπισεν· %
καὶ ἔπεσεν ἐκ τοῦ οὐρανοῦ ἀστὴρ μέγας καιόμενος ὡς λαμπὰς %
καὶ ἔπεσεν ἐπὶ τὸ τρίτον τῶν ποταμῶν καὶ ἐπὶ τὰς πηγὰς τῶν ὑδάτων, %
¹¹ καὶ τὸ ὄνομα τοῦ ἀστέρος λέγεται ὁ \textbf{Ἄψινθος}. %
 καὶ ἐγένετο τὸ τρίτον τῶν ὑδάτων εἰς \textbf{ἄψινθον}, %
 καὶ πολλοὶ τῶν ἀνθρώπων ἀπέθανον ἐκ τῶν ὑδάτων, %
 ὅτι ἐπικράνθησαν. \\[1mm] %
 And the third angel trumpeted. %
 And there fell from the heaven a great star, burning like a lamp, %
 and it fell upon a third of the rivers %
 and upon the springs of water. %
 ¹¹ And the name of the star is called the \textbf{Apsinthos}. %
 And a third of the waters will turn into \textbf{absinthe}, %
 and many of the men will die from the waters, %
 for they were made bitter. \end{quote} % 
\clearpage

In context, the apocalyptic seventh seal had been broken; %
the third of seven trumpet-bearing angels sounds forth, bringing a third set of plagues. %
The context also shares notable features with the [Ps-]Promotus passage. %
Not only do both use the rare second-declension feminine noun, %
but also make clear, proximate references to a lamp, burning, and death. %
If the intertext holds, then John (a late 2nd century!) Revelator (or his redactor) %
borrowed the Promotean image of an absinthe-powered, mosquito-burning contraption %
and turned it into an Exodus-like plague oracle over waters, %
i.e., the places where mosquitoes spread. %
Given the Thracian wordplay in Herodian, %
one also wonders if Revelationʼs Absinthe-star makes a Thracian political/geographical allusion. %
Spartacus, for example, leader of the most famous slave-revolt in Roman history, %
was notably from Thrace (modern Bulgaria). \\[1.5mm] %

\noindent{}\textbf{δυσεντέριον} / \textit{dysenterion} / “dysentery” %
first appears in the TLG, excepting NT texts, %
under the name of Aelius Promotus. %
Dysentery is obviously a common medical term today. %
It occurred frequently among many Greek medical authors. %
Yet before, concurrent with, and even somewhat after Promotus, %
it occurs only as a feminine lemma, %
not a neuter lemma, as Promotus was apparently the first to use. %
In non-Christian literature, the neuter lemma %
next occurs in Moeris (late 2nd or 3rd century), %
who notes the specific distinction between the traditional Attic lemma %
and the new Hellenistic lemma (\textit{Moeridis atticista} δ 34): %

\vfill{}
 \begin{quote}δυσεντερία θηλυκῶς Ἀττικοί· \textbf{δυσεντέριον} Ἕλληνες. \\[1mm] %
 \textit{Dysenteria} feminine Attic [writers]; \textit{\textbf{dysenterion}} Hellenes. %
 \end{quote} % 
 \vfill{}

The mid- to late-2nd century CE Greek medical author Galen %
used only the feminine, not the neuter form. %
He thus remained consistent with the custom of Herodotus, %
Plato, and Hippocrates.\endnote{%
E.g., Herodotus, (\textit{Hist}. 8.115); %
Plato, \textit{Timaeus} 89a; Hippocrates, \textit{passim}.%
} \clearpage %
In the context of Promotusʼ main work, %
Δυναμερόν / \textit{Potency}, %
the neuter gender is not the only distinctive feature %
in his use of the word. %

\vspace*{3mm} 
 \begin{quote}ρκβʹ Πρὸς τεινεσμοὺς καὶ κοιλιακούς --- %
 Προηγεῖται τὸ πλεῖστον \textbf{δυσεντερίων}. Χρὴ δὲ φλεβοτομεῖν εἰ ἐνδέχοιτο, %
 εἰ δὲ μή γε ἀσιτία μέχρι διατρίτου ὠμηλύσεις τε [καὶ] θερμαὶ τιθέμεναι %
 κατ’ ὀσφύος καὶ ἤτρου. Μετὰ δὲ τὰς ἀποκρίσεις ἐγκαθιζέτωσαν εἰς ἀφέψημα %
 τήλεως ἢ μαλάχης ἢ λινοσπέρμου· εἰ δὲ ἀλγοῖεν μετὰ τὰς ἐκκρίσεις,
 ἐγκαθιζέτωσαν εἰς στυμμάτων ἀφέψημα καὶ ἀγαθίδα διὰ στυμμάτων θερμὴν %
 προστιθέσθωσαν <ἢ> πίτυρα ἐν ἀποδέσμῳ δι’ ὀξυκράτου θερμοῦ. (\textit{Potency} 122) \\[1mm] %
 122. For tenesmus and abdominal aches --- %
 Usually this is preceded by \textbf{dysentery}. One should do phlebotomy if possible, %
 otherwise no grain until the third day, applying warm poultices %
 upon the lower back and lower abdomen. Now after the excretions, %
 let them sit in a decoction of fenugreek or mallow or linseed. %
 But if one might feel pain after the secretions, %
 let them sit in a decoction of astringents, % 
 and let them apply a warm skein of astringents <or> %
 husk in a girdle saturated with a warm sour wine mix. \end{quote} % 
 \vspace*{3mm}

Near the end of Acts, when Paul and his sea-faring companions %
are stranded on Malta, %
after he miraculously survives a venomous snakebite, %
Paul is told that the father of Publius, %
the “leading-man” or “chief official” of the island, %
suffered with (neuter) “dysentery” (Acts 28:8--9): % 

\vspace*{3mm} 
 \begin{quote}ἐγένετο δὲ τὸν πατέρα τοῦ Ποπλίου πυρετοῖς καὶ \textbf{δυσεντερίῳ} %
 συνεχόμενον κατακεῖσθαι, πρὸς ὃν ὁ Παῦλος εἰσελθὼν καὶ %
 προσευξάμενος ἐπιθεὶς τὰς χεῖρας αὐτῷ ἰάσατο αὐτόν. %
⁹ τούτου δὲ γενομένου καὶ οἱ λοιποὶ οἱ ἐν τῇ νήσῳ %
ἔχοντες ἀσθενείας προσήρχοντο καὶ ἐθεραπεύοντο. \\[1mm] %
 \textls[-5]{Now it happened that the father of Publius was bed-ridden, %
 occupied with fevers and \textbf{dysentery}. Entering before him, %
 praying, and laying hands upon him, Paul healed him. %
⁹ But after that happened, everyone else on the island %
who had ailments went to him and were cured.} %
\end{quote} % 
 \vspace*{3mm} 

\noindent{}The manuscripts are somewhat divided, %
with the earliest and best witnesses maintaining the neuter form \textit{dusenteriō}, %
yet at least a few dozen minuscules (33 81 104 527 etc) %
reverting to the feminine \textit{dusenteriâ}. %

Manuscripts are important to consider, %
but so is the narrative context and wonder-working depiction in Acts. %
Once bitten by a snake, Paul emerges in a decidedly public, Asclepian role, %
becoming a surpassing, divinely-empowered doctor/priest capable %
of curing all diseases and worthy of the trust of the entire populace. %

If one did not start from the assumption of a 1st or even early 2nd century context for Acts, %
then one would be inclined to find here that author/editor of Acts %
has borrowed (whether directly or indirectly) a rare and perhaps original lemma spelling %
used by the Greek medical author Promotus, %
then used it to show how Paul (likely representing proto-Orthodox Christian priests as miracle-workers more generally) %
surpasses all Greek medical experts and expertise. %
As they often say in operating rooms in the US, %
petty poultices ainʼt got anything on the healing hands of a divine doctor. %